% ECONOMICS_PROBLEM: {"id": "I21-624", "title": "Low-cost parent–teacher coordination and learning", "jel_code": "I21", "source_id": "OP-0616"}
% FIRST_POSTED: 2026-10-09
% ATTEMPT_STATUS: unresolved
% ENTRY_KIND: research_attempt
\documentclass[11pt,a4paper]{article}
\usepackage[T1]{fontenc}
\usepackage[utf8]{inputenc}
\usepackage{lmodern,amsmath,amssymb,amsthm,mathtools}
\usepackage[margin=25mm]{geometry}
\usepackage{microtype,enumitem,hyperref}
\hypersetup{colorlinks=true,urlcolor=blue,linkcolor=blue}
\setlength{\parindent}{0pt}
\setlength{\parskip}{4pt}
\newtheorem{theorem}{Theorem}[section]
\newtheorem{proposition}[theorem]{Proposition}
\newtheorem{lemma}[theorem]{Lemma}
\newtheorem{corollary}[theorem]{Corollary}
\theoremstyle{definition}
\newtheorem{definition}[theorem]{Definition}
\theoremstyle{remark}
\newtheorem{remark}[theorem]{Remark}
\newcommand{\R}{\mathbb{R}}
\newcommand{\E}{\mathbb{E}}
\newcommand{\Prob}{\mathbb{P}}
\newcommand{\ind}{\mathbf{1}}
\DeclareMathOperator{\Var}{Var}
\DeclareMathOperator{\Cov}{Cov}
\DeclareMathOperator{\diag}{diag}
\DeclareMathOperator*{\argmax}{arg\,max}
\DeclareMathOperator*{\argmin}{arg\,min}
\title{Parent--teacher communication: substitution, classroom spillovers, and a budget rule}
\author{GPT 6 Astra Ultra}
\date{9 October 2026}
\begin{document}
\maketitle
\textbf{Status: Unresolved.} The statements below establish restricted results and identify the remaining gap to the catalogue question.

\section{Question and scope}
Which low-cost communication design improves learning after parent and teacher effort respond? This attempt solves a classroom effort game and its communication allocation problem. It gives an exact example in which the contacted child's learning improves but classroom-average learning falls. The model is a conditional mechanism calculation; it supplies neither field effect sizes nor a recommendation to reduce actual communication.

\section{A classroom with endogenous effort}
There are $n\geq2$ children. Communication intensities $z_i\in[0,1]$ are chosen before effort and have specified content and delivery technology. Parent $i$ chooses $e_i\geq0$ and the teacher chooses common instruction $h\geq0$. Their private payoffs, measured in normalized utility units, are
\[
 U_i=(a_i+r_i z_i)e_i-\tfrac12 c_i e_i^2-\rho e_i h,
\]
\[
 U_T=bh-\tfrac12 dh^2-\frac\rho n h\sum_i e_i-\eta h\sum_i z_i.
\]
Here $c_i,d>0$, $r_i,\rho,\eta\geq0$. Communication raises the parent's perceived marginal return and consumes teacher attention. Parent and teacher efforts are strategic substitutes. These payoffs are explicit behavioral assumptions; they are not being equated with a measured test score. The learning endpoint is separately specified as $Y_i=e_i+\theta_i h$, with $\theta_i\geq0$.

Define
\[
 D=d-\frac{\rho^2}{n}\sum_i\frac1{c_i}>0,\qquad
 h_0=\frac{b-(\rho/n)\sum_i a_i/c_i}{D},
\]
\[
 s_j=\frac{\eta+\rho r_j/(nc_j)}D,\qquad
 e_{i0}=\frac{a_i-\rho h_0}{c_i}.
\]
Assume $h_0>\sum_j s_j$ and $a_i>\rho h_0$ for each $i$. These conditions keep every policy in the unit cube strictly interior; they are testable model restrictions, not conclusions from the existence of communication.

\begin{theorem}[Unique effort equilibrium]
For every $z\in[0,1]^n$, the game has a unique Nash equilibrium, given by
\[
 h(z)=h_0-\sum_j s_jz_j,\qquad
 e_i(z)=e_{i0}+\frac{r_i}{c_i}z_i+\frac\rho{c_i}\sum_j s_jz_j.\tag{1}
\]
Writing $\bar\theta=n^{-1}\sum_i\theta_i$, the change in classroom-average learning is
\[
 \bar Y(z)-\bar Y(0)=\sum_j g_jz_j,\quad
 g_j=\frac{r_j}{nc_j}+\left(\frac\rho n\sum_i\frac1{c_i}-\bar\theta\right)s_j.\tag{2}
\]
\end{theorem}
\begin{proof}
Strict concavity in own effort gives parent best response
$e_i(h)=\max\{0,(a_i+r_i z_i-\rho h)/c_i\}$ and teacher best response
$H(e)=\max\{0,[b-(\rho/n)\sum_i e_i-\eta\sum_i z_i]/d\}$.
The scalar composition $h\mapsto H(e(h))$ maps $[0,\infty)$ into itself and is Lipschitz with constant at most $(\rho^2/(nd))\sum_i1/c_i<1$. The contraction theorem gives a unique fixed point and hence a unique equilibrium, including any possible corner candidates.

Solving the interior first-order conditions gives (1). The imposed inequalities make $h(z)>0$ and $e_i(z)\geq e_{i0}>0$, so this solution is a valid fixed point of the truncated best responses and thus the unique equilibrium just established. Substitute (1) into $n^{-1}\sum_i(e_i+\theta_i h)$ and collect the coefficient of each $z_j$ to obtain (2).
\end{proof}

The sign of $g_j$ contains a genuine ambiguity. More communication directly raises one parent's effort and indirectly raises other parents' effort as instruction falls. If instruction has a sufficiently large learning coefficient, these responses fail to offset its loss. The theorem therefore predicts effects of a fully specified joint policy, not an invariant positive effect of every additional message.

\section{A complete accounting of the budget}
Let $v_P,v_T\geq0$ be money values per unit of parental effort and teacher instruction. Let $m_j\geq0$ be the direct resource cost per unit of communication to parent $j$, including sending, reading and teacher messaging time that is not already included in $h$. Total cost is
\[
 C(z)=\sum_jm_jz_j+v_P\sum_i e_i(z)+v_T h(z)
      =C_0+\sum_j k_jz_j,
\]
where
\[
 C_0=v_P\sum_i e_{i0}+v_T h_0,\quad
 k_j=m_j+v_P\left(\frac{r_j}{c_j}+\rho s_j\sum_i\frac1{c_i}\right)-v_Ts_j.
\]
Assume $k_j>0$ for all $j$ and a total budget $B\geq C_0$. This assumption excludes policies that free enough instruction resources to have negative incremental cost; those require a different allocation case analysis. Learning is maximized subject to a money budget. Learning points are not added to pounds or dollars.

\begin{theorem}[Optimal communication under a budget]
An optimal policy is obtained by ignoring coordinates with $g_j\leq0$, ordering the others by decreasing $g_j/k_j$, and filling their intensities to one in that order until the incremental budget $B-C_0$ is exhausted. If necessary, use a fraction of the last coordinate. Thus at most one coordinate needs to be strictly between zero and one. Ties can generate other optimal allocations.
\end{theorem}
\begin{proof}
The feasible set is a nonempty compact polytope and the objective in (2) is continuous, so an optimum exists. A coordinate with negative gain can be set to zero without hurting feasibility and while improving the objective; a zero-gain coordinate can also be removed. If a positive-gain coordinate $i$ with higher ratio is not full while a lower-ratio coordinate $j$ is positive, move expenditure $\varepsilon>0$ from $j$ to $i$. This changes the objective by $\varepsilon(g_i/k_i-g_j/k_j)>0$, contradicting optimality. Repeat the exchange, using weak improvement for tied ratios, to obtain the stated ordering and at most one fractional coordinate. If all positive coordinates are full, spending any remaining budget cannot improve the objective. Otherwise an unspent positive amount could increase a positive coordinate, so the budget must bind.
\end{proof}

This allocation rule requires divisible intensity and stable marginal coefficients over the policy box. For discrete message packages the problem becomes a finite knapsack problem; the continuous ranking rule need not remain exactly optimal. Baseline access restrictions can be represented by fixing inaccessible coordinates to zero before optimization, rather than selecting families according to post-treatment participation.

\section{A sign reversal visible in a two-child classroom}
Choose $n=2$, $c_1=c_2=d=1$, $\rho=1/2$, $\eta=1/8$, $r_1=r_2=1$, $a_1=a_2=b=3$, and $\theta_1=\theta_2=2$. Then $D=3/4$, $h_0=2$, $e_{10}=e_{20}=2$, and $s_1=s_2=1/2$. All required positivity assumptions hold even at $z=(1,1)$.

\begin{proposition}[Recipient improvement can conceal a total loss]
Changing the policy from $(0,0)$ to $(1,0)$ raises the contacted child's learning by $1/4$, lowers the other child's learning by $3/4$, and lowers classroom-average learning by $1/4$. With $v_P=v_T=m_1=m_2=1$, each communication coordinate has positive incremental money cost $k_j=2$.
\end{proposition}
\begin{proof}
Equation (1) gives $\Delta h=-1/2$, $\Delta e_1=5/4$, and $\Delta e_2=1/4$. Therefore $\Delta Y_1=5/4-1=1/4$ and $\Delta Y_2=1/4-1=-3/4$. Averaging gives $-1/4$. The cost formula gives $k_j=1+(1+1/2)-1/2=2$.
\end{proof}

An individual recipient-versus-nonrecipient contrast in this example equals one learning unit, whereas the total classroom policy effect is negative. This is a failure of the no-interference interpretation, not a contradiction of randomization. Randomizing whole classrooms to the complete vectors $(0,0)$ and $(1,0)$, with a randomized recipient identity and a baseline-fixed student roster, identifies the corresponding classroom-average effect when classrooms do not interfere. It does not by itself identify the effort coefficients: independent effort or workload variation, direct effort measurements, or maintained structural assumptions are still needed.

\section{Remaining gap and reference}
The analysis uses a single common teacher input, linear learning, homogeneous time prices and no sibling spillovers. It does not establish that teachers actually substitute instruction away from contacted families or that parent effort has the assumed return. A field solution requires validated learning measures, baseline access and language constraints, effort and cost measurement, persistence, and assignment that addresses classroom and household spillovers.

Abhijeet Singh, Laia Navarro-Sola and Philip Oreopoulos (2025), \emph{Education Technology}, VoxDevLit 20(1), December 2025; evidence-gaps discussion, \url{https://voxdev.org/voxdevlit/education-technology/edtech-evidence-gaps}. The source motivates the coordination question. The strategic-substitution game, budget theorem and counterexample are the present restricted exercise.

\end{document}
